% !TeX spellcheck = es_ES
% !TEX encoding = UTF-8 Unicode
\documentclass[t, compress, spanish, aspectratio=169]{beamer}
%\documentclass[handout]{beamer}    % for the handout

\usetheme[progressbar=frametitle]{metropolis}

%\usetheme[secheader]{Madrid}
%\usetheme{Frankfurt}
%\usetheme{Warsaw}
%\usetheme{Berlin}
%\usetheme{CambridgeUS}

%\usecolortheme{beaver}     % Fondo Blanco y letras Negras - Titulos en Rojo con fondo gris
%\usecolortheme{dove}       % Para el Handout
%\usecolortheme{seahorse}   % Fondo Blanco y letras Negras - Titulos en Celeste
%\usecolortheme{dolphin}    % Fondo Blanco y letras Negras - Titulos en Celeste
%\usecolortheme{seagull}    % Fondo Blanco y letras Negras - Titulos en Gris
%\usecolortheme{crane}      % Fondo Blanco y letras Negras - Titulos en Amarillo
%\usecolortheme{wolverine}  % Fondo Blanco y Letras negras - Titulos en negro con fondo amarillo
%\usecolortheme{lily}       % Fondo Blanco y letras Negras - Titulos en Azul
%\usecolortheme{fly}        % Fondo Gris y letras Negras - Titulos en Rojo
%\usecolortheme{beetle}     % Fondo Gris y letras Negras - Titulos en blanco
%\usecolortheme{albatross}  % Fondo Azul y letras blancas
%\usecolortheme{orchid}     % Default
%\usecolortheme{rose}       % Default
%\usecolortheme{whale}      % Default

\setbeamertemplate{navigation symbols}{} % Gets rid of navigation symbols
%\setbeamercovered{transparent}

\usepackage{dsfont} % Double stroke fonts for math symbols
% \usepackage[T1]{fontenc}  % Helps with hyphenation for languages with accented characters. Note: Does not work with XeLaTex
% \usepackage[latin1]{inputenc} % Allows to type in latin (accents). Note: Does not work with XeLaTex
% \usefonttheme[onlymath]{serif}  % Changes fonts for math to serif
\usepackage[english]{babel} % Manages differences between languages: hyphenation etc

\usepackage{appendixnumberbeamer} % Does not include backup slides in the number of pages

\usepackage{ragged2e}  % Text alignment
\justifying

\usepackage{graphicx} % Format for graphics: additional options for includefigure
\usepackage{pgf} % Tool for graphics

\usepackage{booktabs} % Format for tables: top, bottom and mid lines
\usepackage[scale=2]{ccicons}
\usepackage{threeparttable} % Format for tables
\usepackage{lscape} % Landscape for tables

% Customized commands
  \newcommand{\vs}{\vskip0pt plus.5fill}
  \newcommand{\SR}{\mathbb{R}}
  \DeclareMathOperator{\plim}{plim}
  \DeclareMathOperator{\Var}{Var}
  \DeclareMathOperator{\Avar}{AVar}
  \DeclareMathOperator{\Cov}{Cov}
  \DeclareMathOperator{\E}{\mathbb{E}}
  \newcommand{\pto}{\stackrel{p}{\longrightarrow}}
  \newcommand{\dto}{\stackrel{d}{\longrightarrow}}
  \newcommand{\adistr}{\stackrel{a}{\sim}}
  \newcommand{\epsi}{\varepsilon}
  \newcommand{\tmle}{\hat \theta_{ML}}
  \newcommand{\tgmm}{\hat \theta_{GMM}}
  \newcommand\independent{\protect\mathpalette{\protect\independenT}{\perp}}
  \def\independenT#1#2{\mathrel{\rlap{$#1#2$}\mkern2mu{#1#2}}}



\title [Bootstrap] {Bootstrap}
\author[de Elejalde]{Ramiro de Elejalde}
\institute [UAH-ILADES] {Facultad de Econom\'ia y Finanzas \\ Universidad Alberto Hurtado}
\date[\today] {}
\subject{Empirical IO}
\begin{document}
%------------------------------------------------------------
\begin{frame}
	\titlepage
\end{frame}

% --------------------------------------------------- Slide --
\section{Bootstrap}
\begin{frame}
    \frametitle{Bootstrap (Efron (1979, 1982))}
    \begin{itemize}
    \item <1-> Idea: Obtener muestras aleatorias a partir de la muestra disponible para intentar replicar la aleatoriedad en el muestreo. 
    \item <2-> Para que se utiliza
    \begin{itemize}
        \item Calcular los errores est\'andar cuando es complicado calcularlos con los m\'etodos convencionales.
        \item Obtener refinamientos asint\'oticos.
    \end{itemize}
    \item <3-> Se puede aplicar a una muestra i.i.d. (corte transversal o datos de panel) para estimadores $\sqrt{N}$ consistentes y asintoticamente normales y \emph{suaves}.
    \item <4-> Hay que tener cuidado para aplicarlos a otro tipo de datos (series temporales) y a estimadores que no son suaves (regresi\'on cuart\'ilica, estimaci\'on no param\'etrica, etc.)
    \end{itemize}
\end{frame}
% --------------------------------------------------- Slide --
\subsection{Procedimiento}
\begin{frame}
    \frametitle{Procedimiento}
    \begin{itemize}  [<+->]
        \item $\hat \theta$ es un estimador $\sqrt{N}$ consistente y asint\'oticamente normal de $\theta$.
        \item Muestra i.i.d. $\{y_i,x_i\}_{i=1}^{N}$
        \item Estad\'istico de inter\'es: $s.e.(\hat \theta)$
        \item Procedimiento
        \begin{enumerate}
            \item Tomar $B$ muestras de tama\~no $N$ con reemplazo.
            \item Para cada muestra $b$, estimar $\hat \theta^*_b$.
            \item Calcular 
            \begin{align*}
                \widehat{\Var_{boot}} (\hat \theta)=\frac{1}{B}\sum_{b=1}^B(\hat \theta^*_b-\bar{\hat{\theta}}^*)^2
            \end{align*}
            donde
            \begin{align*}
                \bar{\hat{\theta}}^*= \frac{1}{B}\sum_{b=1}^B \hat \theta^*_b.
            \end{align*}
        \end{enumerate}
    \end{itemize}
\end{frame}
% --------------------------------------------------- Slide --
\subsection{Tipos de bootstrap}
\begin{frame}
    \frametitle{Tipos de bootstrap}
    \begin{itemize} 
        \item <1-> \alert{Bootstrap no param\'etrico}
        \item <2-> \alert{Bootstrap param\'etrico}
        \item[] <2-> Para $y|x \approx F(x,\theta_0)$ obtener muestras de $y|x \sim F(x,\hat \theta)$.
        \item <3-> \alert{Bootstrap residual}
        \item[] <3-> Para el modelo de regresi\'on $y_i=\beta'x_i+u_i$ obtener muestras de $u_i$.
        \item <4-> \alert{Panel bootstrap}
        \item[] <4-> Obtener muestras de individuos (historias completas)
        \item <5-> \alert{Moving blocks}
        \item[] <5-> Para series de tiempo, dividir la muestra en r bloques que no se superponen y obtener muestras de bloques.
        \item <6-> Para calcular el s.e. un n\'umero razonable es 400/500 muestras.
    \end{itemize}
\end{frame}
%% --------------------------------------------------- Slide --
%\subsection{Bootstrap en Stata}
%\begin{frame}
%    \frametitle{Bootstrap en Stata}
%    \begin{itemize}
%        \item Comando  bootstrap \emph{exp}, reps(500) saving(\emph{filename}, replace): \emph{comando}
%        \item Ejemplo: Para la tarea 2 de conducta en la industria de refinamiento de az\'ucar. Calcular
%        \begin{itemize}
%        	\item SE de la media del \'indice de Lerner
%        	\item SE de la estimaci\'on de demanda
%        	\item SE del par\'ametro de conducta
%        \end{itemize}
%        \item Ir a do file en p\'agina web del curso.
%    \end{itemize}
%\end{frame}
\end{document}
